\section{Cohomology of the same actual Cartier cuts}
This slice uses the original base-field-linear Cech complexes and their
actual long exact sequence. It does not redefine Euler characteristic,
cohomology finiteness, or the projective Serre theorem. Every zero scheme
below is the actual zero ideal of the displayed original section.
No cut is assumed reduced or integral.

\begin{definition}[The original section-cut complex]
\label{def:seshadri-original-section-cut-complex}
\lean{EnvelopeProperty.SeshadriAmplenessCriterionCutCohomology.sectionCutSequence}
\leanok
For integral $X$, invertible $A,M$, and $s:\mathcal O_X\to A$ nonzero,
let $I=I(s)$ and $i:I.\mathrm{subscheme}\hookrightarrow X$ be the literal
original zero ideal and inclusion. Define the short complex with terms
$M$, $A\otimes M$, and $i_*i^*(A\otimes M)$ by transporting the actual
closed-ideal sequence tensored with $A\otimes M$. The three whole
comparisons are, in order, the original ideal-to-dual isomorphism followed
by dual evaluation, the actual left unitor, and the existing invertible
projection formula followed by the original unitor on the closed subscheme.
The arrows are the conjugates of the original inclusion and quotient.
\end{definition}
\begin{proof}
The composite of the two conjugated arrows is zero by cancellation of
the middle whole isomorphism and the original inclusion--quotient identity.
All comparisons are constructed from the same original section. This defines
an actual short complex without an incoming exactness or matching certificate.
The final comparison of the right arrow with the standard adjunction-unit
restriction is reserved for the later first-cohomology restriction theorem;
that identity is not an assumption of the higher comparison producer.
\end{proof}

\begin{lemma}[The actual Cartier short exact sequence]
\label{lem:seshadri-actual-cartier-short-exact}
\lean{EnvelopeProperty.SeshadriAmplenessCriterionCutCohomology.sectionCutSequence_shortExact}
\leanok
Let $X$ be integral, let $A,M$ be invertible sheaves, and let
$s:\mathcal O_X\to A$ be nonzero. Set $I=I(s)$, $i:D\hookrightarrow X$.
The original closed-ideal sequence, tensored with $A\otimes M$, and the
canonical evaluation and unit comparisons give an actual short exact sequence
\[
0\longrightarrow M\xrightarrow{\mu_{s,M}} A\otimes M
\xrightarrow{\rho_{s,M}}i_*i^*(A\otimes M)\longrightarrow0.
\]
Here the two arrows are defined by conjugating the actual original
closed-ideal arrows; they are not supplied as certificates.
\uses{def:seshadri-original-section-cut-complex,lem:seshadri-section-cut-euler}
\end{lemma}
\begin{proof}
Use the genuine kernel sequence of
$\mathcal O_X\to i_*\mathcal O_D$. Tensoring with the invertible
$A\otimes M$ is exact. The actual ideal comparison identifies its first
term with $A^{-1}\otimes(A\otimes M)$, and evaluation identifies this
with $M$. The unit identifies the middle term with $A\otimes M$.
The existing invertible projection formula and the actual unit on $D$
identify the quotient with $i_*i^*(A\otimes M)$. Transport the short
complex along these three whole sheaf isomorphisms. The inverse ideal
comparison followed by the original inclusion is evaluation against the
same $s$, so this is the displayed actual multiplication sequence.
\end{proof}

\begin{lemma}[The actual positive cohomology epimorphism]
\label{lem:seshadri-one-cut-cohomology-epi}
\lean{EnvelopeProperty.SeshadriAmplenessCriterionCutCohomology.epi_sectionCut_homologyMap_of_vanishing}
\leanok
\uses{lem:seshadri-actual-cartier-short-exact}
Let $X$ be integral and separated over a field, fix a finite affine open
cover, and use the actual Cartier complex above. If $q\geq1$ and
$H^q(D,i^*(A\otimes M))=0$, its actual left arrow induces an epimorphism
on degree $q$ of the original field-linear Cech cohomology.
\end{lemma}
\begin{proof}
The preimages of the same cover form an affine open cover of $D$, since
the actual closed immersion is affine. Native sheaf-cohomology vanishing
implies exactness of the original Cech complex by the existing affine
comparison. The genuine base-linear comparison and restriction of scalars
preserve its zero homology, and the actual pushforward Cech isomorphism
transports this to the ambient quotient complex. Applying the original
field-complex functor to the quasicoherent Cartier short exact sequence
is short exact. Exactness at its middle degree $q$ group, followed by
the zero quotient group, makes the actual left homology map epic.
\end{proof}

\begin{lemma}[The actual higher cohomology comparison for one cut]
\label{lem:seshadri-one-cut-cohomology-iso}
\lean{EnvelopeProperty.SeshadriAmplenessCriterionCutCohomology.isIso_sectionCut_homologyMap_of_vanishing}
\leanok
Under the Cartier short exact sequence above, suppose $q\geq2$ and
$H^{q-1}(D,i^*(A\otimes M))=H^q(D,i^*(A\otimes M))=0$.
Then the map on degree $q$ original field-linear Cech homology induced
by $\mu_{s,M}$ is an isomorphism.
If only the second vanishing holds, the same map is an epimorphism.
\uses{lem:seshadri-actual-cartier-short-exact,lem:seshadri-one-cut-cohomology-epi}
\end{lemma}
\begin{proof}
The original quasicoherent short exact sequence gives a short exact
sequence of original field-linear Cech complexes on the affine cover.
In its actual long exact sequence, the connecting source in degree
$q-1$ is zero, so the degree $q$ map of the left arrow is monic.
The following group in degree $q$ is zero, so the same map is epic.
The category of modules is balanced. The epimorphism conclusion only
uses exactness at the middle degree $q$ term and its zero following term.
For our use $q\geq2$, both cut vanishings are positive cohomology.
\end{proof}

\begin{lemma}[Two actual Cartier cuts give an arbitrary line difference]
\label{lem:seshadri-two-cut-euler-difference}
\lean{EnvelopeProperty.SeshadriAmplenessCriterion.eulerCharacteristic_cartierDifference}
\uses{lem:seshadri-section-first-difference}
\leanok
Let $X$ be an integral projective scheme over a field, and let $L,B,M$
be invertible sheaves. Suppose $a\ne0$ is an actual section of $L\otimes B$
and $b\ne0$ is an actual section of $B$, with their actual zero schemes
$D_a,D_b$. Then
\[
 \chi(X,L\otimes M)-\chi(X,M)
 =\chi(D_a,(L\otimes M)|_{D_a})-\chi(D_b,M|_{D_b}).
\]
\end{lemma}
\begin{proof}
Set $F=B^{-1}\otimes M$. The actual section exact sequences for $a$
and $b$, tensored with this same $F$, respectively give
\[
 \chi(X,(L\otimes B)\otimes F)-\chi(X,F)
 =\chi(D_a,((L\otimes B)\otimes F)|_{D_a})
\]
and
\[
 \chi(X,B\otimes F)-\chi(X,F)
 =\chi(D_b,(B\otimes F)|_{D_b}).
\]
The actual evaluation of $B$ and associativity identify the two positive
middle terms with $L\otimes M$ and $M$, respectively. Their pullbacks
provide the same identifications on the original closed schemes.
Subtracting the second equality from the first proves the assertion.
\end{proof}

\begin{lemma}[The original arbitrary line powers have actual cut differences]
\label{lem:seshadri-arbitrary-power-cut-difference}
\lean{EnvelopeProperty.SeshadriAmplenessCriterion.exists_eulerCharacteristic_powerCartierDifference}
\uses{lem:seshadri-two-cut-euler-difference}
\leanok
Let $X$ be an integral projective scheme over a field, and let $L$ be any
original invertible sheaf. There exist an actual invertible $B$, a nonzero
section $a$ of $L\otimes B$, and a nonzero section $b$ of $B$, such that,
for every $n\geq0$, with the same actual zero schemes $D_a,D_b$,
\[
 \chi(X,L^{n+1})-\chi(X,L^n)
 =\chi(D_a,L^{n+1}|_{D_a})-\chi(D_b,L^n|_{D_b}).
\]
\end{lemma}
\begin{proof}
Use an actual projective closed embedding. Ordinary coherent global
generation supplies a sufficiently large nonnegative projective twist
$B$ for which the original $L\otimes B$ is globally generated. The
original coordinate monomials generate this same $B$. On a nonempty
scheme a generated invertible sheaf has a nonzero member of its actual
generating family, since otherwise its epimorphic evaluation and hence
its identity would vanish. Choose such members $a,b$. These constructions
are the existing original nonzero-generator producers, not additional
assumptions. Apply the preceding equality to $M=L^n$ for each $n$ and use
the original tensor-power successor isomorphism, also on its actual
pullback to $D_a$. Integrality makes both zero ideals effective Cartier;
neither zero scheme is assumed integral or reduced.
\end{proof}

\begin{lemma}[Two actual cuts compare the same original line twists]
\label{lem:seshadri-two-cut-cohomology-iso}
\lean{EnvelopeProperty.SeshadriAmplenessCriterionCutCohomology.twoCutCohomologyIso}
\leanok
Let $L,B,M$ be invertible on an integral separated scheme over $k$,
$a:\mathcal O_X\to L\otimes B$ and $b:\mathcal O_X\to B$ nonzero.
Write $i:D_a\hookrightarrow X$ and $j:D_b\hookrightarrow X$ for their
actual zero subschemes. Suppose $q\geq2$ and the four original groups
\[
\begin{gathered}
H^{q-1}(D_a,i^*(L\otimes M)),\quad H^q(D_a,i^*(L\otimes M)),\\
H^{q-1}(D_b,j^*M),\quad H^q(D_b,j^*M)
\end{gathered}
\]
vanish. Then the original degree $q$ field-linear cohomology of
$L\otimes M$ and $M$ are isomorphic. The isomorphism is the actual span
through $F=B^{-1}\otimes M$, formed from the maps of the same two cuts.
\uses{lem:seshadri-one-cut-cohomology-iso,lem:seshadri-two-cut-euler-difference}
\end{lemma}
\begin{proof}
Use the two actual Cartier sequences with the same $F$.
Associativity and evaluation of $B$ give the canonical whole comparisons
$(L\otimes B)\otimes F\cong L\otimes M$ and $B\otimes F\cong M$,
and their pullbacks give the corresponding actual cut comparisons.
Both left arrows induce isomorphisms in degree $q$ by the preceding lemma.
Compose the inverse of the first with the second and these whole
comparisons. This construction does not assert an $H^1$ comparison:
two surjections with a common source do not yield a map between their targets.
\end{proof}

\begin{lemma}[Eventual higher comparisons from genuine cut ampleness]
\label{lem:seshadri-eventual-two-cut-cohomology-iso}
\lean{EnvelopeProperty.SeshadriAmplenessCriterionCutCohomology.exists_twoCut_powerCohomologyIso_of_ample}
\leanok
In the preceding setting suppose the actual restrictions $i^*L$ and
$j^*L$ are ample in the native positive-power-embedding sense. There is
one integer $N$ such that, for every $n\geq N$ and every $q\geq2$,
the degree $q$ original field-linear groups of $L^{n+1}$ and $L^n$
are isomorphic by the preceding actual two-cut span. The sections
$B,a,b$ and their cuts remain fixed before all $n,q$.
\uses{lem:seshadri-two-cut-cohomology-iso}
\end{lemma}
\begin{proof}
Unpack the native ample datum on each actual cut into its positive
power, original closed projective embedding, and whole original
$\mathcal O(1)$ comparison. In the existing namespace
\texttt{OriginalAmpleSerreVanishing}, apply its theorem
\texttt{exists\_sheafCohomology\_bound} with the actual unit sheaf. This theorem already handles the finitely
many power residues internally and gives one bound for every positive
cohomology degree. Remove the unit by its actual unitor and transport
along the existing canonical pullback-of-original-powers isomorphism.
Taking the maximum of the two bounds makes the four positive groups
in the preceding lemma vanish, including the $n+1$ restriction on $D_a$.
The actual original power-successor isomorphism completes the claim.
Ampleness of $L$ on $X$ is not an assumption.
\end{proof}

\begin{definition}[The same section's original power complex]
\label{def:seshadri-same-section-original-power-complex}
\lean{EnvelopeProperty.SeshadriAmplenessCriterionCutCohomology.powerSectionSequence}
\leanok
Given an invertible $L$, a nonnegative integer $r$, and the same nonzero
section $s:\mathcal O_X\to L^r$, define the $n$-th power complex to have
terms $L^n$, $L^{n+r}$, and $i_*i^*L^{n+r}$, where $i$ is its original
zero-subscheme inclusion. Transport the original section-cut complex
for $A=L^r$ and $M=L^n$ by the identity on its first term, the canonical
symmetry and tensor-power addition isomorphism on its second, and the
pushforward of the pullback of that same whole comparison on its third.
\end{definition}
\begin{proof}
The existing tensor-power addition isomorphism uses the recursively formed
original powers. The functorial pullback and pushforward therefore preserve
both the actual power sheaf and the literal same zero subscheme. Transport
of the original short-complex zero relation defines the two displayed arrows.
\end{proof}

\begin{lemma}[Exactness of the same original power complex]
\label{lem:seshadri-same-section-original-power-short-exact}
\lean{EnvelopeProperty.SeshadriAmplenessCriterionCutCohomology.powerSectionSequence_shortExact}
\leanok
\uses{def:seshadri-same-section-original-power-complex,lem:seshadri-actual-cartier-short-exact}
Every original power complex just defined is short exact.
\end{lemma}
\begin{proof}
Both original powers are invertible. The actual Cartier section sequence
is short exact, and the three internally constructed whole comparisons
define an isomorphism of short complexes. Transport its short exactness.
\end{proof}

\begin{lemma}[The genuine single-cut eventual first cohomology epimorphisms]
\label{lem:seshadri-single-cut-eventual-hone-epi}
\lean{EnvelopeProperty.SeshadriAmplenessCriterionCutCohomology.exists_singleCut_powerHOne_epi_of_ample}
\leanok
Let $L$ be invertible, let $r$ be nonnegative, and let $s:\mathcal O_X\to L^r$ be a
nonzero original section. For its same actual zero subscheme
$i:D\hookrightarrow X$, suppose $i^*L$ is natively ample. There is
$N$ such that, for every $n\geq N$, the actual multiplication map of
this one section induces an epimorphism
\[
H^1(X,L^n)\longrightarrow H^1(X,L^{n+r})
\]
on the original field-linear cohomology.
\uses{lem:seshadri-one-cut-cohomology-iso}
\end{lemma}
\begin{proof}
Apply the actual Cartier sequence with $A=L^r$ and $M=L^n$.
The canonical power-addition isomorphism identifies its middle term
with $L^{r+n}=L^{n+r}$. The same native ample Serre producer on $D$
makes $H^1(D,i^*L^{n+r})$ zero uniformly for sufficiently large $n$.
The epimorphism part of the one-cut lemma applies. No two-cut $H^1$
comparison is used, and the existence of $s$ is left to the later
Euler-growth argument that actually constructs it.
\end{proof}

\begin{lemma}[Finite residue stabilization and eventual restriction surjectivity]
\label{lem:seshadri-single-cut-hone-stable-restriction-surjective}
In the preceding setting suppose $r>0$ and $X$ is projective over $k$. There is
$N'$ such that, for every $n\geq N'$, the same $H^1$ multiplication map
is an isomorphism, and the actual restriction map
\[
H^0(X,L^{n+r})\longrightarrow H^0(D,i^*L^{n+r})
\]
is surjective.
\uses{lem:seshadri-single-cut-eventual-hone-epi}
\end{lemma}
\begin{proof}
The original global sections identify naturally with the kernel of the first
Cech differential by sheaf gluing. The incoming differential in degree zero
is zero, so the canonical cycles-to-homology isomorphism identifies this
same kernel with original $H^0$. Naturality follows by composing the actual
augmentation with each product projection: both composites are the same
restriction of the original section map. On an invertible affine frame,
the transported quotient and the standard pullback--pushforward adjunction
unit both take a coefficient to its residue modulo the same zero ideal.
These actual local equalities glue, providing the standard restriction
comparison internally. No comparison certificate is supplied.
Genuine projective coherent cohomology finiteness gives finite-dimensional
original field-linear $H^1$ groups. On each of the finitely many residue
classes modulo $r$, the dimensions are eventually nonincreasing natural
numbers, by the actual eventual epimorphisms. Their attained least value
makes each sequence eventually constant. Take one bound for all residues.
A surjective map of equal finite dimensions is injective, hence each
same actual $H^1$ map is now an isomorphism. In the actual long exact
sequence its preceding connecting map from $H^0(D,i^*L^{n+r})$ is zero,
since its composite with this monomorphism is zero. Exactness makes
the original $H^0$ restriction epimorphic, hence surjective on sections.
\end{proof}
